%% parts/p11/philosophy-of-logic.tex --- Part 11
\chapter{Philosophy of Logic}\label{sec:p11-philosophy-of-logic}

Philosophical disputes about logic often turn on a change of object: truth in a
fixed structure, validity over a class of structures, derivability in a
calculus, rational revision, or linguistic practice. This chapter states each
object before comparing positions. A label such as ``classical,'' ``pluralist,''
or ``deflationary'' never substitutes for the formal commitments it is meant
to summarize.

\section{Truth, metalanguages, and fixed points}
\label{sec:p11-truth}

\begin{hbdefinition}{def:p11-tarski-truth}{Tarskian truth definition and Convention T}
Let \(L\) be a fully interpreted object language and \(M\) a metalanguage able
to name its expressions and describe its semantics. A materially adequate
truth definition entails every instance
\[
 \ulcorner\varphi\urcorner\text{ is true in }L
 \quad\longleftrightarrow\quad
 \varphi
\]
for sentences \(\varphi\) of \(L\); this adequacy demand is Convention T. The
object/metalanguage distinction is necessary because allowing the truth
predicate and its semantic resources without restriction in one language
recreates the liar construction. The definition stands as a metalinguistic
scheme, not an object-language axiom schema, differs from a compositional
recursive truth definition, and is used to state Tarski's hierarchy and the
undefinability boundary \cite{Tarski1983,Halbach2011}.
\index{Convention T}\index{truth!Tarskian}\index{object language}\index{metalanguage}
\end{hbdefinition}

For a fixed first-order structure, the compositional clauses define truth by
recursion on formula complexity. For sufficiently strong arithmetical
languages, no arithmetical formula defines the set of all true sentences of
that same language. This undefinability theorem does not prohibit truth in a
stronger metalanguage or restricted truth predicates for lower complexity
classes.

\begin{hbdefinition}{def:p11-kripke-truth}{Kripke fixed-point truth}
Fix a bivalent base structure, extend its language by a partial unary truth
predicate, and evaluate sentences by the strong-Kleene compositional clauses.
A truth state \(v\) supplies disjoint positive and negative extensions of
that predicate; the base interpretation remains fixed. The truth jump
\(\Gamma(v)\) computes the new predicate extensions from sentence values.
A fixed point satisfies \(\Gamma(v)=v\); transfinite iteration starts with
the truth predicate undefined on the whole domain, not with every ordinary
predicate undefined. A sentence is \emph{grounded} when it receives a
classical value at that least fixed point. These names are needed to separate
partial semantic evaluation, fixed-point choice, and groundedness. They stand
as semantic constructions, differ from a Tarskian language hierarchy by
allowing self-reference with truth-value gaps, and are used to diagnose liar-
like sentences \cite{Kripke1975,Martin1984}.
\index{truth!Kripke fixed point}\index{grounded sentence}\index{strong Kleene scheme}
\end{hbdefinition}

The least fixed point is not the only fixed point, and different evaluation
schemes need not have the same fixed points. Kripke's supervaluation schemes
evaluate over the specified bivalent extensions of a partial interpretation,
not merely over its fixed points. Their consistency and eligibility conditions
are given in Section~\ref{sec:p11-kripke-fixed-points}. A maximal fixed
point, a supervaluation jump and a revision sequence answer different
questions. Revision theories iterate a classical hypothesis about
the truth predicate and study stable, nearly stable, and periodic behaviour;
they do not identify a single least partial model.

\begin{proofboundary}{def:p11-axiomatic-truth}{Axiomatic truth-theory comparison boundary}
The literature abbreviations \(\mathsf{TB}\), \(\mathsf{CT}\),
\(\mathsf{KF}\), and \(\mathsf{FS}\) do not name unique formal theories until
the arithmetical base, expanded language, coding, induction scope, and full
axiom/rule package are fixed. This chapter therefore uses them only as
retrieval labels for four families:
\begin{center}
\begin{tabularx}{\textwidth}{@{}l X X@{}}
\toprule
Family & Distinguishing commitment & Variation that must be fixed for a theorem \\
\midrule
TB-style & sentence-level disquotational biconditionals & base theory, sentence class, and induction in the truth language \\
CT-style & uniform compositional clauses for syntax constructors & coding, full constructor list, induction, and whether compositionality applies to truth-language formulas \\
KF-style & positive/negative strong-Kleene fixed-point axioms & coding guards, complete dual clauses, induction, and fixed-point/minimality principles \\
FS-style & compositional axioms plus truth-introduction and truth-elimination rules & rule domain, compositional package, base, and induction \\
\bottomrule
\end{tabularx}
\end{center}
The distinction is necessary because variants in one row can have different
conservativity and proof-theoretic strength. These are metalinguistic survey
families, not four locally axiomatized theories; they are nonredundant through
the displayed commitments and are used only in the following comparison. Exact
formal variants and their strength results are catalogued in
\cite{Halbach2011}.
\index{axiomatic truth theory}\index{compositional truth}\index{fixed-point truth}
\end{proofboundary}

Disquotational biconditionals may be conservative over a base under restricted
induction, while compositional truth with stronger induction need not be; the
exact result is variant-specific and remains cited rather than asserted for a
bare acronym. Accordingly, ``deflationary'' does not by itself determine a
formal theory or its strength.

\section{Semantic paradoxes}
\label{sec:p11-paradoxes}

\begin{center}
\begin{tabularx}{\textwidth}{@{}l X X@{}}
\toprule
Paradox & Formal mechanism & What a response must change \\
\midrule
liar & diagonal sentence \(L\leftrightarrow\neg T(\ulcorner L\urcorner)\) & truth language, semantic values, or inference principles \\
Curry & \(C\leftrightarrow(T(\ulcorner C\urcorner)\to\bot)\) & contraction/conditional reasoning or truth principles \\
Yablo & sentence \(Y_n\) says every later \(Y_m\) is false & global infinite dependency, not a displayed self-reference \\
Grelling & predicates applied to their own linguistic names & type/level discipline for semantic predicates \\
Berry & definability plus a size-bound description & exact syntax, coding, and definability predicate \\
\bottomrule
\end{tabularx}
\end{center}

The strengthened liar can be formulated so that common gap replies yield a
new revenge sentence. Curry's paradox does not use negation and therefore
tests structural and conditional principles rather than only bivalence
\cite{Curry1942}. Yablo's sequence avoids a sentence referring directly to
itself, but its infinitary dependency still requires a semantic totality
\cite{Yablo1993}. Grelling and Berry require careful distinctions among
expressions, their names, semantic predicates, and effective length bounds.

There is no single ``solution strategy'' without a target theory. A hierarchy
restricts syntax; fixed-point semantics permits gaps; supervaluationism
quantifies over admissible precisifications; paraconsistent theories reject
explosion; revision theories replace a static valuation by a dynamics; and
axiomatic truth theories restrict principles inside arithmetic. Each response
must state which classical laws, truth clauses, structural rules, or
metalinguistic resources it preserves \cite{Priest2006,
BeallGlanzbergRipley2018}.

\section{Vagueness and the sorites}
\label{sec:p11-vagueness}

\begin{hbdefinition}{def:p11-sorites}{Sorites series}
A sorites series fixes objects \(a_0,\ldots,a_n\), a vague predicate \(P\),
clear endpoint judgments \(P(a_0)\) and \(\neg P(a_n)\), and tolerance premises
\(P(a_i)\to P(a_{i+1})\). The name is needed because the contradiction depends
on the entire chain of local tolerance claims, not on vagueness alone. It
stands as an argument schema, differs from a single borderline case, and is
used to compare the theories below.
\index{sorites paradox}\index{vagueness}
\end{hbdefinition}

Supervaluationism considers a set of admissible classical precisifications.
A sentence is supertrue when true on all of them and superfalse when false on
all. This preserves classical tautologies at the global supervaluation level,
but local consequence at a single precisification and global consequence over
all admissible precisifications must not be conflated \cite{Fine1975,
Keefe2000}. Degree theories assign values in an ordered algebra and must state
the chosen conditional and consequence threshold; a fuzzy value is not a
probability of an unknown sharp value \cite{Hajek1998,Smith2008}.

Epistemicism retains a sharp boundary while claiming that its location is
unknowable, so it changes the epistemology rather than the classical extension
\cite{Williamson1994}. Contextual accounts allow the comparison class or
standard to shift with conversational context. No view follows merely from
the existence of borderline cases; each must validate its own treatment of
tolerance, higher-order vagueness, and logical consequence.

\section{Logical consequence and logical constants}
\label{sec:p11-consequence}

\begin{hbdefinition}{def:p11-model-consequence}{Model-theoretic consequence}
Relative to a signature and class \(\mathcal K\) of interpretations,
\(\Gamma\models_{\mathcal K}\varphi\) means that every member of
\(\mathcal K\) satisfying every sentence in \(\Gamma\) satisfies \(\varphi\).
The subscript is necessary because varying models, admissible reinterpretations,
or the domain policy changes the relation. It stands as a semantic relation
among sentences, differs from derivability in a calculus, and is used to state
soundness, completeness, and objections to model-theoretic analyses
\cite{Tarski1983,Etchemendy1990}.
\index{logical consequence!model-theoretic}
\end{hbdefinition}

A proof-theoretic account instead explains consequence through canonical
introduction and elimination rules. \emph{Harmony} requires, at minimum, that
elimination rules extract no more from a connective than its introduction
rules warrant, often made precise by local reduction and expansion properties.
Harmony is therefore a relation among specified rules, not a synonym for
proof-theoretic validity \cite{Dummett1991,Prawitz1965}.

Permutation invariance offers a semantic demarcation of logical constants:
operations invariant under all permutations of a domain do not distinguish
particular objects. Generalized quantifiers require corresponding invariance
conditions across bijections of structures. Invariance explains one formal
sense of topic neutrality, but choosing the transformation group and semantic
type is part of the proposal \cite{Sher1991}.

Logical monism holds that one consequence relation is correct for the target
purpose; pluralism allows more than one equally legitimate relation, usually
because distinct cases preserve distinct features. Relativism, instrumental
choice, and the technical fact that many consequence relations can be defined
are weaker claims. A pluralist position must identify each relation, its cases,
and its norm of correctness \cite{BeallRestall2006}.

\section{Second-order logic, plurality, and commitment}
\label{sec:p11-ontology}

\begin{hbdefinition}{def:p11-plural-quantification}{Plural quantification}
Plural quantification adds variables interpreted plurally and locutions such
as ``\(x\) is one of the \(xx\)'' without first treating \(xx\) as a single
set-valued object. The terminology is needed to distinguish many-object
reference from singular quantification over a collection. It stands as a
formal and semantic proposal, differs from standard set-valued second-order
semantics in its claimed ontology, and is used to assess whether full
second-order quantification commits one to sets \cite{Boolos1984,Shapiro1991}.
\index{plural quantification}\index{second-order logic!ontological status}
\end{hbdefinition}

Under full semantics, second-order variables range over every subset or
relation on the first-order domain; under Henkin semantics, their ranges are
specified components of a many-sorted structure. These logics have different
validities and metatheorems. Calling full second-order logic ``set theory in
disguise'' emphasizes its strong semantic resources, whereas a plural reading
denies that plural reference requires a single set object. Neither slogan
changes the formal distinction.

Quine's criterion reads existentially quantified variable values as the
ontological commitments of a regimented theory \cite{Quine1948}. Applying it
requires a translation, a background logic, and a criterion of theoretical
equivalence. Nominalism, structuralism, and platonism then differ over what
mathematical quantification requires; they do not follow solely from the
completeness or incompleteness of a calculus.

Identity likewise separates formal and metaphysical questions. First-order
logic with identity fixes a congruence-like semantic relation interpreted as
actual equality in each structure. Relative identity, counterpart relations,
and intensional identity principles require different languages and cannot be
silently substituted for that clause.

\section{Methodological boundary}
\label{sec:p11-method-boundary}

The positions surveyed here are organized by explicit formal choices:

\begin{itemize}
  \item object language, metalanguage, and semantic resources;
  \item truth values or partial valuation scheme;
  \item consequence relation and structural rules;
  \item quantifier range and interpretation policy;
  \item preserved classical principles and intended application.
\end{itemize}

A philosophical label earns no theorem status. Historical attribution is
cited; formal claims are checked against the displayed definitions; and
normative conclusions remain arguments rather than consequences of notation.
